\documentclass[10pt,a4paper]{amsart}
\usepackage[margin=19mm]{geometry}
\usepackage{amsmath,amssymb,mathtools}
\usepackage[scr=rsfs,cal=euler]{mathalfa}
\usepackage{xfrac}
\usepackage[unicode,hidelinks,bookmarks=false]{hyperref}
\newcommand{\ie}{i.e.\ }
\newcommand{\cf}{cf.\ }
\newcommand{\resp}{respectively\ }
\newcommand{\Subsection}{\subsection}
\providecommand{\nobreakdash}{\nobreak\hbox{-}}
\setlength{\emergencystretch}{2em}
\multlinegap0pt
\usepackage{amssymb}
\usepackage{xcolor}
\usepackage{mathtools}
\newtheorem{lem}{Lemma}
\newcommand{\eqdef}{\coloneqq}
\newcommand{\id}{\operatorname{Id}}
\title{From Score Matching to Diffusion: \\ A Fine-Grained Error Analysis in the Gaussian Setting}
\author{Samuel Hurault, Matthieu Terris, Thomas Moreau, Gabriel Peyr\'e}
\date{Source: arXiv:2503.11615v3 (CC BY 4.0)}
\begin{document}
\maketitle

\noindent\emph{Source and licence.} From Score Matching to Diffusion: \\ A Fine-Grained Error Analysis in the Gaussian Setting. Version arXiv:2503.11615v3 (CC BY 4.0), distributed by arXiv under the Creative Commons Attribution 4.0 International licence, http://creativecommons.org/licenses/by/4.0/, as recorded in the arXiv licence metadata for this exact version and shown on the abstract page https://arxiv.org/abs/2503.11615v3. Full licence notice: \texttt{licenses/arxiv-2503.11615-CC-BY-4.0.txt}.\par\smallskip

\noindent\textit{Editorial scope.} Counted unit: the article's lem lem:Langevin\_1 (source0.tex lines 1393--1410). The article's complete original proof of it is delivered with it (source0.tex lines 1411--1512, one \texttt{proof} environment of 102 lines). No mathematics is written, repaired or re-derived: the statement and the proof are the article's own lines, in the article's own notation, and no retained line is substituted or re-flowed.  Retained uncounted as the source-local context the proof needs: lines 1318--1320.  Nothing was omitted from any retained span, so the package carries no omission record.  No retained line contains a citation, so the wrapper carries no bibliography and no bibliography entry.  No retained line contains a figure environment, an image-inclusion command or drawing code, so no image is delivered and the package declares no asset dependency.  Editorial formatting only, in addition: an AMS \texttt{amsart} preamble that restates the article's own declarations and macros used by the retained lines, namely the environments \texttt{lem} on separate counters, together with the standard packages those lines need.  The retained environments print in this document as Lemma 1 in order of appearance, whereas the article prints the same items with the numbering of its own full document.  The complete native source of the article as distributed in the arXiv e-print for this exact version is bundled unchanged as \texttt{sources/174773/source0.tex}.
\begin{align} \label{eq:inv_lyap}
    L^\gamma_C[X] = CX + XC^\top - \gamma C X C^\top .
\end{align}

\begin{lem} \label{lem:Langevin_1}
With the previous notations, the mean $\mu_\varepsilon$ and the covariance $\Sigma_\varepsilon$ expand as 
\begin{align*}
\mu_\varepsilon &=  \mu + \varepsilon C_\sigma \delta   + o(\varepsilon)  \\
\text{and} \quad\Sigma^\gamma_\varepsilon &= \Sigma_0^\gamma  + \varepsilon \Sigma^\gamma_1(\Delta) + \varepsilon^2 \Sigma^\gamma_2(\Delta) + o(\varepsilon^2),
\end{align*}
with 
 \begin{align} \label{eq:Sigma_0}
\Sigma_0^\gamma &= \Sigma^\gamma_{C_\sigma^{-1}} = C_\sigma\left( \id - \frac{\gamma}{2} C_\sigma^{-1} \right)^{-1} \\
\Sigma_1^\gamma(\Delta) &= -\Bigl(L^\gamma_{C_\sigma^{-1}}\Bigr)^{-1}\Bigl[
\Delta \Sigma^\gamma_{C_\sigma^{-1}} (\id - \gamma {C_\sigma^{-1}}) + (\id - \gamma {C_\sigma^{-1}})\Sigma^\gamma_{{C_\sigma^{-1}}} \Delta^\top
\Bigr] \label{eq:Sigma_1}
\\
\Sigma_2^\gamma(\Delta) &= -\Bigl(L^\gamma_{C_\sigma^{-1}}\Bigr)^{-1}\Bigl[
\Delta \Sigma^\gamma_{{1}}(\Delta) (\id - \gamma {C_\sigma^{-1}}) + (\id - \gamma {C_\sigma^{-1}})\Sigma^\gamma_{{1}}(\Delta) \Delta^\top- \gamma   \Delta  \Sigma^\gamma_{{C_\sigma^{-1}}}  \Delta^\top
\Bigr]. \label{eq:Sigma_2}
\end{align}
\end{lem}

\begin{proof}
We start with the mean. 
 \begin{align*}
 \mu_\varepsilon &\eqdef (A_\varepsilon)^{-1}b_\varepsilon = \left( A^* + \varepsilon \Delta + o(\varepsilon) \right)^{-1} (b^* + \varepsilon\Delta\mu + \varepsilon \delta + o(\varepsilon)) \\
 &= (A^*)^{-1}b^* - \varepsilon (A^*)^{-1} \Delta (A^*)^{-1} b^* + \varepsilon (A^*)^{-1} \delta  + \varepsilon (A^*)^{-1} \Delta \mu + o(\varepsilon^2)
 \end{align*}
Using $A^* = C_\sigma^{-1}$ and $b^* = C_\sigma^{-1} \mu$, we get 
  \begin{align*}
 \mu_\varepsilon &= \mu + \varepsilon C_\sigma  \delta  + o(\varepsilon)
 \end{align*}
We now focus on the covariance.
 \begin{align*}
\Sigma^\gamma_{A} = \Bigl(L^\gamma_{A^*+\varepsilon\Delta}\Bigr)^{-1}[2I],
\end{align*}
Substitute $A = A^* + \varepsilon\Delta$ into the definition~\eqref{eq:inv_lyap} of $L^\gamma_A[Z]$:
 \begin{align*}
L^\gamma_{A^*+\varepsilon\Delta}[Z] = (A^*+\varepsilon\Delta)Z + Z(A^*+\varepsilon\Delta)^\top - \gamma  (A^*+\varepsilon\Delta)Z (A^*+\varepsilon\Delta)^\top.
 \end{align*}
We now expand each term:
 \begin{align*}
\begin{aligned}
(A^*+\varepsilon\Delta)Z &= A^* Z + \varepsilon   \Delta Z,\\
Z(A^*+\varepsilon\Delta)^\top &= Z (A^*)^\top + \varepsilon   Z \Delta^\top,\\
(A^*+\varepsilon\Delta) Z (A^*+\varepsilon\Delta)^\top 
&= A^*Z (A^*)^\top + \varepsilon\Bigl[\Delta Z (A^*)^\top + A^*Z \Delta^\top\Bigr] + \varepsilon^2   \Delta Z \Delta^\top.
\end{aligned}
 \end{align*}
Substituting these into the expression for $L^\gamma_{A^*+\varepsilon\Delta}[Z]$ yields
 \begin{align*}
L^\gamma_{A^*+\varepsilon\Delta}[Z] &= A^*Z + Z (A^*)^\top - \gamma   A^*Z (A^*)^\top\\ \nonumber
&\quad {}+ \varepsilon\Bigl[\Delta Z + Z \Delta^\top - \gamma\Bigl(\Delta Z (A^*)^\top + A^*Z \Delta^\top\Bigr)\Bigr] - \gamma  \varepsilon^2   \Delta Z \Delta^\top.
 \end{align*}
For clarity, we define the following operators:
 \begin{align*}
\mathcal{L}_0[Z] &\eqdef A^*Z + Z (A^*)^\top - \gamma   A^*Z (A^*)^\top,\\
\mathcal{L}_1[Z] &\eqdef \Delta Z + Z \Delta^\top - \gamma\Bigl(\Delta Z (A^*)^\top + A^*Z \Delta^\top\Bigr),\\
\mathcal{L}_2[Z] &\eqdef - \gamma   \Delta Z \Delta^\top.
 \end{align*}
such that
 \begin{align*}
L^\gamma_{A^*+\varepsilon\Delta}[Z] = \mathcal{L}_0[Z] + \varepsilon  \mathcal{L}_1[Z] + \varepsilon^2  \mathcal{L}_2[Z].
 \end{align*}
We now seek an expansion of the solution:
 \begin{align*}
\Sigma^\gamma_{A^*+\varepsilon\Delta} = \Sigma_0^\gamma + \varepsilon  \Sigma_1^\gamma + \varepsilon^2  \Sigma_2^\gamma + o(\varepsilon^2).
 \end{align*}
By definition,
 \begin{align*}
L^\gamma_{A^*+\varepsilon\Delta}\Bigl[\Sigma^\gamma_{A^*+\varepsilon\Delta}\Bigr] = 2I.
 \end{align*}
Substituting the expansions of both the operator and the solution, we obtain:
 \begin{align*}
\Bigl(\mathcal{L}_0 + \varepsilon  \mathcal{L}_1 + \varepsilon^2  \mathcal{L}_2\Bigr)
\Bigl[\Sigma_0^\gamma + \varepsilon  \Sigma_1^\gamma + \varepsilon^2  \Sigma_2^\gamma\Bigr] = 2I.
 \end{align*}
We now collect terms by powers of $\varepsilon$.

\paragraph{Order \(\varepsilon^0\)}
 \begin{align*}
\mathcal{L}_0\Bigl[\Sigma_0^\gamma\Bigr] = 2I.
 \end{align*}

Thus, the zeroth-order term is given by
 \begin{align*}
\Sigma_0^\gamma = \Bigl(L^\gamma_{A^*}\Bigr)^{-1}[2I].
 \end{align*}

\paragraph{Order \(\varepsilon^1\)}
Collecting the terms of order $\varepsilon$, we have:
 \begin{align*}
\mathcal{L}_0\Bigl[\Sigma_1^\gamma\Bigr] + \mathcal{L}_1\Bigl[\Sigma_0^\gamma\Bigr] = 0.
 \end{align*}
Therefore,
 \begin{align*}
\Sigma_1^\gamma = -\Bigl(L^\gamma_{A^*}\Bigr)^{-1}\Bigl[\mathcal{L}_1\Bigl[\Sigma_0^\gamma\Bigr]\Bigr],
 \end{align*}
or explicitly,
 \begin{align*}
\Sigma_1^\gamma = -\Bigl(L^\gamma_{A^*}\Bigr)^{-1}\Bigl[
  \Delta  \Sigma_0^\gamma + \Sigma_0^\gamma  \Delta^\top - \gamma\Bigl(\Delta  \Sigma_0^\gamma  (A^*)^\top + A^*  \Sigma_0^\gamma  \Delta^\top\Bigr)
\Bigr].
 \end{align*}

\paragraph{Order \(\varepsilon^2\)}
At order $\varepsilon^2$, we have:
 \begin{align*}
\mathcal{L}_0\Bigl[\Sigma_2^\gamma\Bigr] + \mathcal{L}_1\Bigl[\Sigma_1^\gamma\Bigr] + \mathcal{L}_2\Bigl[\Sigma_0^\gamma\Bigr] = 0.
 \end{align*}
This implies
 \begin{align*}
\Sigma_2^\gamma = -\Bigl(L^\gamma_{A^*}\Bigr)^{-1}\Bigl[
\mathcal{L}_1\Bigl[\Sigma_1^\gamma\Bigr] + \mathcal{L}_2\Bigl[\Sigma_0^\gamma\Bigr]
\Bigr],
 \end{align*}
or in an expanded form,
 \begin{align*}
\Sigma_2^\gamma = -\Bigl(L^\gamma_{A^*}\Bigr)^{-1}\Bigl[
  \Delta  \Sigma_1^\gamma + \Sigma_1^\gamma  \Delta^\top - \gamma\Bigl(\Delta  \Sigma_1^\gamma  (A^*)^\top + A^*  \Sigma_1^\gamma  \Delta^\top\Bigr)
- \gamma   \Delta  \Sigma_0^\gamma  \Delta^\top
\Bigr].
 \end{align*}
\end{proof}

\end{document}
